\scanpage{27}
\deleted{\unsure{Corollaire} $A$ \illegible{2}}

\uline{Prop. 6} Soit $A\in\underline b^+$, on a
\[
\begin{aligned}
(5)\qquad \varphi_A^-(t)
 &=\inf_{\substack{E\in P(\mathcal A)\\ \operatorname{Tr}E<t}}
       \sup_{x\in(1-E)H}(Ax,x)\\
 &=\inf_{\substack{E\in P(\mathcal A)\\ \operatorname{Tr}E<t}}
       \|(1-E)A(1-E)\|.
\end{aligned}
\]

\uline{Démonstration} Soit $\underline C$ \inserted{$=C_0(M)$} une sous-algèbre commutative auto-adjointe maximale dans $\underline b$, contenant $A$. Si on se borne dans \unsure{le calcul de son} $\operatorname{Inf}$ aux $E$ \unsure{pris} dans $\underline C$, on trouve
\[
\inf_{\substack{\unsure{E\mathrm{\ comm.\ }A}\\ \operatorname{Tr}(E)<t}}
\|(1-E)A\|_\infty\geq\varphi_A^-(t)
\]
\margin{$\unsure{\Downarrow}$}
(N° 1, formule (4)) d'où \unsure{donc} dans (5, $\varphi_A^-(t)$ est $\geq$ le \unsure{deuxième} mb.

\uline{Corollaire 1} $A\in\underline b$ implique
\margin{--- pour $|A|=\sqrt{A^*A}$}
\[
\begin{aligned}
(6)\qquad \varphi_{|A|}^-(t)
 &=\inf_{\substack{E\in P(\mathcal A)\\ \operatorname{Tr}(E)<t}}
       \sup_{x\in(1-E)H}\|Ax\|\\
 &=\inf_{\substack{E\in P(\mathcal A)\\ T(E)<t}}\|A(1-E)\|.
\end{aligned}
\]

\uline{Corollaire 2} $A\in\underline b$, $B\in\mathcal A$ impliquent
\[
(7)\qquad
\left\{\begin{aligned}
\varphi_{|BA|}&\leq\|B\|\,\varphi_{|A|}\\
\varphi_{|AB|}&\leq\|B\|\,\varphi_{|A|}.
\end{aligned}\right.
\]

\uline{Corollaire 3} $\varphi_A$ est fonction croissante de $A\in\underline b^+$.

\uline{Corollaire 4} Les fonctions $A\ \text{---}\ \varphi_{|A|}^-(t)$ sur $\underline b$ forment un ensemble équicontinu de fonctions (quand $\underline b$ est \unsure{munie} de sa topologie \unsure{uniforme}).
\margin{$\Big|$}

\margin{\uline{Corollaire 5}\quad $\varphi_{|1+A|}\leq\varphi_{1+|A|}$}
\mspage{28}
